\section{Migrar la receta de escalado de LLM a modalidades múltiples}

En los dos capítulos anteriores se establecieron las interfaces; este capítulo discute la ruta de entrenamiento. La clase no consiste en inventar desde cero un conjunto de ciencia multimodal, sino de desglosar el éxito de los LLM en cuatro hipótesis verificables: si el seguimiento de instrucciones puede cruzar modalidades, si el razonamiento puede aprovechar contexto visual/auditivo, si el aumento de escala mejora continuamente y si una arquitectura dispersa puede controlar los costos. Cada una podría ser cierta, pero también podría fallar debido a la alta dimensionalidad, redundancia y diferencias de objetivos en las modalidades.

\subsection{Los cuatro pilares de migración}

Esta sección avanza la interfaz de secuencia unificada del capítulo anterior hacia la ruta de entrenamiento: primero se aclara qué experiencia de los LLM merece ser migrada y luego se verifican las condiciones de fallo posibles para cada una en modalidades de alta dimensión. Al leer la figura, debe entenderse por capas de capacidades conductuales, efectos de escala e infraestructura, en lugar de considerar las cuatro columnas como una lista de tareas completadas.

\lecturefigure{slide-019.jpg}{Los cuatro principios de escalado multimodal derivados del paradigma de LLM}{Diapositiva 19 del paquete oficial; clase 00:09:16--00:11:33}

\begin{knowledgebox}{Lectura de la figura: leer en orden desde "conducta" hasta "infraestructura"}
La primera columna trata sobre prompting y seguimiento de instrucciones, resolviendo cómo el usuario controla al modelo; la segunda es planificación multimodal, que resuelve cómo el modelo organiza pasos intermedios entre diferentes fuentes de información; la tercera aborda escalado de datos/modelos, planteando si el rendimiento mejora continuamente con la escala; la cuarta es escalado arquitectónico, usando MoE, dispersión de atención o especialización de modalidades para reducir los cómputos activos. Las dos primeras columnas modifican principalmente el comportamiento, mientras que las dos últimas determinan si el entrenamiento y la inferencia son viables.
\end{knowledgebox}

\begin{importantbox}{Hipótesis de escalado, no garantía de escalado}
Si aumentar los datos o los parámetros genera una curva de rendimiento $Q(N,D)$, la experiencia con modelos de lenguaje sugiere que $\partial Q/\partial N>0$ y $\partial Q/\partial D>0$ se cumplen en un intervalo. Sin embargo, la calidad de los datos multimodales, errores de sincronización, fotogramas repetidos, anotaciones faltantes y el *mismatch* de *loss* pueden alterar la curva. Lo verdaderamente importante por verificar es el escalado con cómputo emparejado (*matched-compute scaling*), no solo mostrar las puntuaciones finales de modelos más grandes.
\end{importantbox}

\teachervoice{00:09:16--00:11:33, el docente considera que seguimiento de instrucciones, planificación, escala y dispersión arquitectónica son principios migrados desde los LLM. El tono más importante aquí es "enfoque": estas son rutas de investigación, no una lista de ingeniería ya resuelta.}

\subsection{Dos ejes de problemas organizan la exposición completa}

El capítulo anterior enumeró las recetas de escalado; este capítulo las sintetiza en dos preguntas verificables: cómo se generan salidas no textuales y si el entrenamiento conjunto produce transferencia de capacidades. Las diferencias entre modelos posteriores pueden compararse volviendo a estos dos ejes; los lectores también deben registrar por separado la representación, el *loss* y el compartimiento de parámetros, evitando que el concepto de "modelo unificado" oculte las distintas condiciones experimentales.

\lecturefigure{slide-021.jpg}{Las dos preguntas de investigación centrales del modelo Omni}{Diapositiva 21 del paquete oficial; clase 00:11:37--00:12:18}

\begin{knowledgebox}{Lectura de la figura: generación y transferencia son dos problemas independientes}
La primera pregunta es "cómo generar modalidades no textuales", centrándose en la representación y el *loss*; la segunda pregunta si la comprensión y la generación tienen una transferencia positiva, centrándose en si el backbone compartido genera sinergia. Un modelo puede resolver la primera pregunta pero presentar transferencia negativa en la segunda. Chameleon y Transfusion exploran principalmente la primera; MoT maneja simultáneamente los conflictos de capacidad; la segunda se discute frontalmente solo en la segunda mitad de la conferencia.
\end{knowledgebox}

A continuación, el diseño del modelo puede escribirse como un quintuple:
$$
\mathcal{M}=(R, O, L, P, C),
$$
donde $R$ es representación, $O$ es ordenamiento/interfaz de secuencia, $L$ son pérdidas, $P$ es compartimiento/ruteo de parámetros y $C$ es asignación de cómputo y capacidad. Los diferentes artículos no se sustituyen mutuamente, sino que hacen elecciones distintas en las cinco dimensiones.

\begin{knowledgebox}{Leer la tabla unificada de preguntas para los artículos posteriores}
Al ver un nuevo modelo multimodal, primero pregunte: ¿cuáles son las modalidades de entrada y salida? ¿La imagen es código discreto o latente continuo? ¿Es autoregresivo entre modalidades? ¿Qué pérdida se usa para cada salida? ¿Qué atención/FFN se comparte? ¿El ruteo es determinista o aprendido? Al comparar, ¿los parámetros activos, los tokens, los FLOPs y los datos son coincidentes? Si no hay respuestas a estas preguntas, mirar solo la frase "modelo unificado" no tiene suficiente significado técnico.
\end{knowledgebox}

\subsection{Resumen de la sección}

El objetivo del escalado multimodal no es forzar todas las señales en el mismo vocabulario, sino preservar la interfaz de control unificada de los modelos de secuencia, al tiempo que permite representaciones específicas de modalidad, objetivos y capacidad dispersa. Las tres familias de modelos siguientes responden: ¿hasta dónde puede llegar la discretización completa? ¿Cómo coexisten los objetivos de generación continua con los modelos de lenguaje? ¿Cuándo debe el compartimiento de parámetros ser reemplazado por dispersión determinista?
